% ============================================================
% 第3章 化学平衡（对应课件：选必一2.2 化学平衡、2.3 化学反应方向、2.4 化学反应调控）
% ============================================================
\chapter{化学平衡}

上一章讨论了反应进行的快慢，本章讨论可逆反应进行的限度。先学会判断反应是否达到平衡，再用平衡常数 $K$、$K_p$ 和转化率描述反应进行的程度，然后分析条件改变后平衡怎样移动、怎样读平衡图像，接着用 $\Delta G=\Delta H-T\Delta S$ 判断反应能否自发进行，最后讨论工业生产中反应条件的选择和多个反应同时发生的多重平衡。

广东卷第19题主要考本章内容：2025年考了 $\Delta G$ 判断、平衡组成图、$K$ 的计算和催化剂选择性，2026年考了吸收平衡的 $K$ 计算，2024年第17题考了“如何证明平衡已建立”。选择题常考平衡标志、平衡移动和图像。本章分值约 10～14 分。下一章的电离平衡、水解平衡和沉淀溶解平衡，都要用到本章的平衡常数和平衡移动。

\section{化学平衡状态}

\subsection{可逆反应与平衡的建立}
在同一条件下，既能向正反应方向进行、又能向逆反应方向进行的反应叫可逆反应，用“\ce{<=>}”表示。可逆反应不能进行到底，任一时刻反应物与生成物共存。注意 \ce{2H2 + O2 ->[\text{点燃}] 2H2O} 与 \ce{2H2O ->[\text{电解}] 2H2 ^ + O2 ^} 条件不同，不是可逆反应。

以 \ce{N2 + 3H2 <=> 2NH3} 为例：开始时反应物浓度最大，$\vz$ 最大，$\vn=0$；随着反应进行，$\vz$ 减小、$\vn$ 增大，直到 $\vz=\vn$，各物质浓度不再改变，此时反应达到\textbf{化学平衡状态}。

\begin{center}
\begin{tikzpicture}[scale=0.75,font=\small]
\begin{scope}
\draw[-Stealth] (0,0)--(5,0) node[below left]{$t$}; \draw[-Stealth] (0,0)--(0,3.2) node[left]{$c$};
\draw[thick,mainblue] (0,2.8) .. controls (1,1.8) and (1.8,1.6) .. (2.6,1.6)--(4.8,1.6) node[right]{\ce{H2}};
\draw[thick,deepgreen] (0,1.0) .. controls (1,0.7) and (1.8,0.62) .. (2.6,0.62)--(4.8,0.62) node[right]{\ce{N2}};
\draw[thick,warnred] (0,0) .. controls (1,0.6) and (1.8,0.75) .. (2.6,0.76)--(4.8,0.76) node[above right]{\ce{NH3}};
\draw[dashed] (2.6,0)--(2.6,3) node[above]{$t_1$};
\end{scope}
\begin{scope}[xshift=7cm]
\draw[-Stealth] (0,0)--(5,0) node[below left]{$t$}; \draw[-Stealth] (0,0)--(0,3.2) node[left]{$v$};
\draw[thick,mainblue] (0.02,2.8) .. controls (1,1.7) and (1.8,1.4) .. (2.6,1.4)--(4.8,1.4);
\draw[thick,warnred] (0.02,0.02) .. controls (1,1.0) and (1.8,1.38) .. (2.6,1.4);
\node[mainblue] at (0.9,2.4) {$\vz$}; \node[warnred] at (1.1,0.5) {$\vn$};
\node at (3.8,1.75) {$\vz=\vn\neq0$};
\draw[dashed] (2.6,0)--(2.6,3) node[above]{$t_1$};
\end{scope}
\end{tikzpicture}
\end{center}

\begin{gainian}[化学平衡状态的特征：逆、等、动、定、变]
“逆”指研究对象是可逆反应；“等”指 $\vz=\vn\neq0$，这是平衡的本质；“动”指平衡是动态的，反应并没有停止；“定”指各组分的浓度（含量）保持不变，这是平衡的外在表现；“变”指条件改变时平衡会被破坏而发生移动。化学平衡是该条件下反应的最大限度，决定了反应物的最大（平衡）转化率。
\end{gainian}

\subsection{平衡状态的判断}
判断依据只有两条：一是“正逆相等”，二是“变量不变”。

\textbf{(1) 正逆相等：}同一物质的生成速率等于消耗速率；用不同物质表示时，必须一个表示正反应、一个表示逆反应，且速率之比等于计量数之比。以 \ce{N2 + 3H2 <=> 2NH3} 为例：“$2\vz(\ce{N2})=\vn(\ce{NH3})$”是平衡标志；“$3v(\ce{N2})=v(\ce{H2})$”任何时刻都成立，不能作标志；“断裂 3 mol H—H 键的同时断裂 6 mol N—H 键”表示正逆速率相等，是平衡标志；“断裂 3 mol H—H 键的同时形成 6 mol N—H 键”都表示正反应，不是标志。

\textbf{(2) 变量不变：}一个量若随反应进行而变化（变量），则它不变时反应达到平衡；若它本来就不随反应变化（定量），则不能作为标志。

\begin{center}
\small
\begin{tabularx}{\linewidth}{p{3.4cm}XXX}
\toprule
物理量（不变时） & \ce{N2(g) + 3H2(g) <=> 2NH3(g)}（气体分子数变化） & \ce{H2(g) + I2(g) <=> 2HI(g)}（气体分子数不变） & \ce{C(s) + H2O(g) <=> CO(g) + H2(g)}（有固体） \\
\midrule
恒温恒容：总压强 $p$ & 是 & 否（$n$ 总不变） & 是 \\
恒温恒压：总体积 $V$ & 是 & 否 & 是 \\
气体总物质的量 $n$ & 是 & 否 & 是 \\
恒容：气体密度 $\rho=m_\text{气}/V$ & 否（$m$、$V$ 都不变） & 否 & 是（气体质量变化） \\
平均摩尔质量 $\bar M=m_\text{气}/n_\text{气}$ & 是 & 否 & 是 \\
有色物质的颜色 & — & 是（\ce{I2} 浓度变化） & — \\
\bottomrule
\end{tabularx}
\end{center}

\begin{yicuo}
\begin{enumerate}
\item “各物质浓度之比等于计量数之比”“各物质浓度相等”都不一定是平衡状态。平衡时各物质浓度保持不变，但不一定相等，也不一定成比例。
\item 反应物为固体、生成的气体组成比例固定的分解反应要特别小心：如 \ce{NH4Cl(s) <=> NH3(g) + HCl(g)}，两种气体始终按 1:1 生成，$\bar M$ 和 \ce{NH3} 的体积分数始终不变，不能作标志；但恒容时压强、气体密度是变量，可以作标志。
\item 绝热容器中，反应热使体系温度变化，“温度不变”可作为平衡标志；同时平衡常数也会随温度变化。
\end{enumerate}
\end{yicuo}

\begin{liti}[平衡标志综合判断]
可逆反应 \ce{2NO2(g) <=> 2NO(g) + O2(g)} 在恒容密闭容器中进行：①单位时间内生成 $n$ mol \ce{O2} 的同时生成 $2n$ mol \ce{NO2}；②单位时间内生成 $n$ mol \ce{O2} 的同时生成 $2n$ mol NO；③用 \ce{NO2}、NO、\ce{O2} 表示的速率之比为 2:2:1；④混合气体颜色不再改变；⑤混合气体密度不再改变；⑥混合气体压强不再改变；⑦混合气体平均相对分子质量不再改变。能说明达到平衡状态的是\tk。
\tcblower
\jie ①生成 \ce{O2} 为正反应、生成 \ce{NO2} 为逆反应，且比例 1:2 正确，是。②都是正反应，否。③任何时刻都成立，否。④\ce{NO2} 红棕色，颜色不变即浓度不变，是。⑤恒容且全为气体，密度始终不变，否。⑥气体分子数变化，压强是变量，是。⑦$m$ 不变、$n$ 变化，$\bar M$ 是变量，是。答案：①④⑥⑦。
\end{liti}

\begin{liti}[恒温恒压下的体积计算]
某体积可变的密闭容器中盛有适量 A 和 B 的混合气体，在一定条件下发生 \ce{A(g) + 3B(g) <=> 2C(g)}。维持温度和压强不变，达到平衡时容器体积为 $V$ L，其中 C 的体积分数为 10\%。则原混合气体的体积为\tk L，平衡时 A 消耗了\tk L。
\tcblower
\silu 恒温恒压下，气体体积之比等于物质的量之比，可直接用“体积”列三段式。\\
\jie 平衡时 $V(\mathrm{C})=0.1V$。每生成 2 体积 C，气体总体积减少 $1+3-2=2$ 体积，所以减少 $0.1V$；原混合气体为 $V+0.1V=1.1V$ L。消耗 A 为 $0.1V\times\frac12=0.05V$ L，消耗 B 为 $0.15V$ L。\\
\dianbo “差量法”：总体积的变化量与某物质变化量成比例，可省去设未知数。
\end{liti}

\section{化学平衡常数}

\subsection{平衡常数的表达式}
800 ℃ 时，对 \ce{CO(g) + H2O(g) <=> CO2(g) + H2(g)}，无论起始投料如何，平衡时 $\dfrac{c(\ce{CO2})\cdot c(\ce{H2})}{c(\ce{CO})\cdot c(\ce{H2O})}$ 都约等于 1.0。大量实验表明这一规律具有普遍性。

\begin{gainian}[化学平衡常数]
一定温度下，可逆反应 \ce{$m$A(g) + $n$B(g) <=> $p$C(g) + $q$D(g)} 达到平衡时，生成物浓度幂之积与反应物浓度幂之积的比值是一个常数：
\[K=\frac{c^{\,p}(\mathrm{C})\cdot c^{\,q}(\mathrm{D})}{c^{\,m}(\mathrm{A})\cdot c^{\,n}(\mathrm{B})}\qquad(\text{各浓度均为平衡浓度})\]
对指定反应，$K$ \textbf{只与温度有关}，与浓度、压强、催化剂无关。
\end{gainian}

书写 $K$ 的表达式时要注意：
\begin{enumerate}
\item \textbf{固体、纯液体不写入表达式}（其“浓度”视为常数）。如 \ce{CaCO3(s) <=> CaO(s) + CO2(g)}，$K=c(\ce{CO2})$。
\item 稀水溶液中的反应，水不写入；非水体系中有水参与时（如酯化反应 \ce{CH3COOH + C2H5OH <=> CH3COOC2H5 + H2O}），水要写入。
\item 离子反应先写成离子方程式再写 $K$。如 \ce{Cl2 + H2O <=> H+ + Cl- + HClO}，$K=\dfrac{c(\ce{H+})c(\ce{Cl-})c(\ce{HClO})}{c(\ce{Cl2})}$。
\item $K$ 与方程式的写法一一对应：
\end{enumerate}
\begin{center}
\small
\begin{tabular}{ll}
\toprule
方程式的变化 & $K$ 的变化 \\
\midrule
计量数扩大 $n$ 倍 & $K'=K^n$ \\
正、逆反应 & $K_\text{逆}=1/K_\text{正}$ \\
两式相加（相减） & $K=K_1\cdot K_2$（$K=K_1/K_2$） \\
\bottomrule
\end{tabular}
\end{center}
这可以和盖斯定律对照记忆：方程式相加，$\Delta H$ 相加，$K$ 相乘。

\subsection{平衡常数的意义与应用}
①衡量反应进行的程度。$K$ 越大，正反应进行得越完全，反应物平衡转化率越大。一般 $K>10^5$ 认为反应基本完全。如 25 ℃ 时 \ce{H2 + F2 <=> 2HF} 的 $K\approx6.5\times10^{95}$，几乎完全；\ce{H2 + I2 <=> 2HI} 的 $K\approx867$，明显不完全。但 $K$ 只反映“限度”（热力学），不反映快慢（动力学）：$K$ 很大的反应也可能很慢。

②判断反应方向。任意时刻的浓度幂之积比值叫\textbf{浓度商} $Q$（形式同 $K$，但用任意时刻浓度）。
\[Q<K\ \text{反应正向进行};\qquad Q=K\ \text{平衡状态};\qquad Q>K\ \text{反应逆向进行}\]

③判断反应的热效应。升高温度，$K$ 增大，则正反应吸热；$K$ 减小，则正反应放热。

④计算平衡组成和转化率。

\begin{gainian}[转化率、产率与选择性]
\[\alpha(\text{反应物})=\frac{\text{转化的量}}{\text{起始的量}}\times100\%,\qquad \text{产率}=\frac{\text{实际产量}}{\text{理论产量}}\times100\%\]
\[\text{目标产物的选择性}=\frac{\text{转化为目标产物的反应物的量}}{\text{反应物转化的总量}}\times100\%,\qquad \text{目标产物产率}=\alpha\times\text{选择性}\]
\end{gainian}

\begin{liti}[三段式与 $Q$、$K$]
将 0.1 mol CO 与 0.1 mol \ce{H2O(g)} 充入 10 L 密闭容器，800 ℃ 达平衡时 $c(\ce{CO})=0.005$ mol·L$^{-1}$。(1) 求 $K$ 和 CO 的平衡转化率。(2) 同温下，CO 仍为 0.1 mol，\ce{H2O(g)} 改为 0.3 mol，求 CO 的平衡转化率。(3) 同温下，CO 为 0.1 mol、\ce{H2O(g)} 为 0.375 mol，某时刻测得 $c(\ce{CO})=0.0025$ mol·L$^{-1}$，此时是否平衡？反应向哪个方向进行？
\tcblower
\jie (1) 起始 $c(\ce{CO})=c(\ce{H2O})=0.01$，转化 0.005，平衡时四种物质均为 0.005 mol·L$^{-1}$，$K=\dfrac{0.005\times0.005}{0.005\times0.005}=1$，$\alpha(\ce{CO})=50\%$。\\
(2) 设 CO 转化 $x$ mol·L$^{-1}$：$\dfrac{x^2}{(0.01-x)(0.03-x)}=1$，解得 $x=0.0075$，$\alpha(\ce{CO})=75\%$。增大 \ce{H2O} 的浓度，CO 的转化率提高（而 \ce{H2O} 自身的转化率由 50\% 降为 25\%）。\\
(3) 此时 CO 转化了 0.0075，$c(\ce{H2O})=0.0375-0.0075=0.03$，$Q=\dfrac{0.0075^2}{0.0025\times0.03}=0.75<K=1$，未平衡，反应正向进行。
\end{liti}

\begin{liti}[平衡常数的组合与反应方向]
已知：①\ce{2H2(g) + CO(g) <=> CH3OH(g)} 的 $K_1$ 在 500 ℃、800 ℃ 时分别为 2.5、0.15；②\ce{CO2(g) + H2(g) <=> CO(g) + H2O(g)} 的 $K_2$ 分别为 1.0、2.5。(1) 写出 ③\ce{CO2(g) + 3H2(g) <=> CH3OH(g) + H2O(g)} 的 $K_3$ 与 $K_1$、$K_2$ 的关系，并判断反应③的 $\Delta H$ 正负。(2) 500 ℃ 时，测得反应③某时刻 \ce{H2}、\ce{CO2}、\ce{CH3OH}、\ce{H2O} 的浓度（mol·L$^{-1}$）分别为 0.8、0.1、0.3、0.15，判断此时 $\vz$ 与 $\vn$ 的大小。
\tcblower
\jie (1) ③$=$①$+$②，$K_3=K_1\cdot K_2$。500 ℃：$K_3=2.5$；800 ℃：$K_3=0.375$。升温 $K_3$ 减小，正反应放热，$\Delta H<0$。\\
(2) $Q=\dfrac{0.3\times0.15}{0.1\times0.8^3}=\dfrac{0.045}{0.0512}\approx0.88<K_3=2.5$，反应正向进行，$\vz>\vn$。
\end{liti}

\subsection{压强平衡常数 $K_p$}
对气相反应，常用平衡分压代替平衡浓度，得到\textbf{压强平衡常数} $K_p$。分压 $=$ 总压 $\times$ 物质的量分数，即 $p(\mathrm{B})=p_\text{总}\cdot\dfrac{n(\mathrm{B})}{n_\text{总}}$。如 \ce{N2(g) + 3H2(g) <=> 2NH3(g)}：
\[K_p=\frac{p^2(\ce{NH3})}{p(\ce{N2})\cdot p^3(\ce{H2})}\]
$K_p$ 同样只与温度有关。有些题目把分压除以标准压强 $p^\ominus=100$ kPa 使其无单位（如 2026 年广东卷第19题），按题给定义计算即可。

\begin{kanwu}[课件 2.2 第38～39页]
课件中压强平衡常数、物质的量分数平衡常数对应的方程式写作“\ce{$m$A(g) + $n$B(g) <=> $e$C(g) + $f$d(g)}”，其中“d”应为“D”。课件给出的关系 $K_p=K_c(RT)^{(e+f)-(m+n)}$ 与 $K_x=K_p\cdot p_\text{总}^{(m+n)-(e+f)}$ 正确，了解即可，高考一般直接用分压计算 $K_p$。
\end{kanwu}

\begin{fangfa}[求 $K_p$ 的步骤]
①列三段式求平衡时各气体的物质的量和气体总物质的量（别漏掉不反应的稀释气体如 Ar、水蒸气）；②求各物质的量分数；③乘以总压得分压；④代入表达式。若反应前后气体分子数相等，总压和 $n_\text{总}$ 可约去，直接用物质的量代入。
\end{fangfa}

\begin{liti}[$K_p$ 的计算]
一定温度、恒定总压 $p=100$ kPa 下，\ce{N2O4(g) <=> 2NO2(g)} 达到平衡时 \ce{N2O4} 的转化率为 50\%。求 $K_p$；若保持温度不变，将总压改为 200 kPa，平衡转化率变大还是变小？
\tcblower
\jie 设起始 \ce{N2O4} 1 mol，平衡时 \ce{N2O4} 0.5 mol、\ce{NO2} 1 mol，$n_\text{总}=1.5$ mol。$p(\ce{N2O4})=\frac13p$，$p(\ce{NO2})=\frac23p$，$K_p=\dfrac{(\frac23p)^2}{\frac13p}=\frac43p=133$ kPa（一般式：$K_p=\dfrac{4\alpha^2}{1-\alpha^2}p$）。总压增大，平衡向气体分子数减少的方向（逆向）移动，转化率变小；由 $\dfrac{4\alpha^2}{1-\alpha^2}\times200=133$ 可解得 $\alpha\approx0.38$。
\end{liti}

\section{化学平衡的移动}

\subsection{平衡移动的本质}
条件改变使 $\vz\neq\vn$，原平衡被破坏，在新条件下建立新平衡的过程叫\textbf{平衡移动}。$\vz>\vn$ 时正向移动，$\vz<\vn$ 时逆向移动。从热力学看，就是条件改变使 $Q\neq K$：$Q<K$ 正向移动，$Q>K$ 逆向移动。用“$Q$ 与 $K$”分析，往往比“速率”更清楚。

\subsection{浓度、压强、温度、催化剂的影响}
\textbf{(1) 浓度：}增大反应物浓度或减小生成物浓度，$Q$ 减小，$Q<K$，平衡正向移动；反之逆向移动。实验：\ce{Fe^3+ + 3SCN- <=> Fe(SCN)3}（红色），加 KSCN 溶液红色加深；加铁粉（\ce{2Fe^3+ + Fe = 3Fe^2+}，消耗 \ce{Fe^3+}）红色变浅。注意：加 KCl 固体，\ce{K+}、\ce{Cl-} 不参与反应，平衡不移动；改变固体或纯液体的量，平衡也不移动。

\textbf{(2) 压强：}只影响有气体参与、且反应前后气体分子数不等的反应。对 \ce{2NO2(g) <=> N2O4(g)}，压缩体积到原来的一半，各浓度都变为 2 倍：
\[Q=\frac{2c(\ce{N2O4})}{[2c(\ce{NO2})]^2}=\frac{K}{2}<K\]
平衡正向（气体分子数减少的方向）移动。现象：颜色先变深（\ce{NO2} 浓度瞬间增大）后变浅（平衡正移），但最终仍比原来深（只能“减弱”不能“抵消”）。对 \ce{H2(g) + I2(g) <=> 2HI(g)}，压缩时 $Q=K$ 不变，平衡不移动，但颜色因浓度增大而加深。

\textbf{(3) 温度：}升高温度，平衡向吸热方向移动；降低温度，向放热方向移动。温度是\textbf{唯一能改变 $K$} 的因素。实验：\ce{2NO2(g) <=> N2O4(g)} $\Delta H<0$，热水中红棕色加深，冰水中变浅。

\textbf{(4) 催化剂：}同等程度改变 $\vz$、$\vn$，平衡不移动，但缩短达到平衡的时间。

\begin{gainian}[勒夏特列原理（平衡移动原理）]
如果改变影响平衡的一个因素（浓度、压强、温度），平衡就向着能够\textbf{减弱}这种改变的方向移动。适用于已达平衡的体系、只改变一个条件的情况；移动的结果只能减弱、不能抵消这种改变。
\end{gainian}

\begin{benzhi}[“充入惰性气体”平衡怎样移动]
\textbf{恒温恒容}：总压增大，但各反应物、生成物的浓度（分压）不变，$Q=K$，平衡不移动。\textbf{恒温恒压}：容器体积增大，各气体浓度（分压）同比例减小，相当于减压，平衡向气体分子数增大的方向移动。工业上常在恒压下通入水蒸气（如乙苯脱氢制苯乙烯）或 Ar 稀释，就是为了提高气体分子数增大的反应的平衡转化率。
\end{benzhi}

\subsection{平衡移动与转化率的变化规律}
\begin{center}
\small
\begin{tabularx}{\linewidth}{p{5.6cm}X}
\toprule
操作 & 结果 \\
\midrule
增大反应物 A 的浓度（反应 \ce{A + B <=> C}） & 平衡正移，B 的转化率增大，A 自身的转化率减小 \\
恒温恒容，按原比例同时增加所有反应物 & 相当于加压：若正反应气体分子数减小，各反应物转化率增大；分子数不变，转化率不变 \\
恒温恒容，再加入一定量的唯一反应物（如 \ce{2NO2 <=> N2O4} 中再加 \ce{NO2}） & 相当于加压：\ce{NO2} 的转化率增大（新平衡中 \ce{N2O4} 的体积分数增大） \\
恒温恒压，按原比例投料 & 新平衡与原平衡等效，各组分的百分含量和转化率不变 \\
\bottomrule
\end{tabularx}
\end{center}

\begin{fangfa}[“虚拟容器法”（放缩法）]
恒容容器中按比例增加投料，可设想先把新增物质放在另一个相同容器中，在相同条件下建立与原来完全相同的平衡，再把两个容器“压”到一个容器中，这样“加料”问题就变成了“加压”问题。同理，把一个容器的体积扩大，可看作“减压”。用同样的思路可以理解\textbf{等效平衡}：恒温恒压下，投料比例相同即等效；恒温恒容下，气体分子数不变的反应按比例投料即等效（百分含量相同），分子数改变的反应需“一边倒”后投料量完全相同才等效。
\end{fangfa}

\begin{liti}[等效平衡]
恒温恒容的甲、乙、丙三个相同容器中发生 \ce{2SO2(g) + O2(g) <=> 2SO3(g)}：甲中投入 2 mol \ce{SO2} 和 1 mol \ce{O2}，乙中投入 2 mol \ce{SO3}，丙中投入 4 mol \ce{SO2} 和 2 mol \ce{O2}。比较：(1) 甲、乙平衡时 \ce{SO3} 的体积分数；(2) 甲中 \ce{SO2} 的转化率 $\alpha_\text{甲}$ 与乙中 \ce{SO3} 的转化率 $\alpha_\text{乙}$ 之和；(3) 甲、丙中 \ce{SO2} 的转化率。
\tcblower
\jie (1) 把乙中 2 mol \ce{SO3} “一边倒”为 2 mol \ce{SO2} 和 1 mol \ce{O2}，与甲完全相同，恒温恒容下两者为等效平衡，平衡组成完全相同，\ce{SO3} 体积分数相等。(2) 设甲平衡时 \ce{SO3} 为 $x$ mol，则 $\alpha_\text{甲}=\frac x2$；乙平衡时 \ce{SO3} 也为 $x$ mol，$\alpha_\text{乙}=\frac{2-x}{2}$，二者之和为 1。(3) 丙相当于把两个甲“压缩”到一个容器中，加压平衡正移，$\alpha_\text{丙}>\alpha_\text{甲}$。
\end{liti}

\begin{liti}[浓度变化倍数的分析]
某温度下，密闭容器中发生 \ce{$a$X(g) + $b$Y(g) <=> $c$Z(g) + $d$W(g)}，达平衡后将容器体积扩大为原来的 2 倍，再次达到平衡时 W 的浓度为原平衡的 0.65 倍。则平衡向\tk{}方向移动，$a+b$\tk$c+d$，X 的转化率\tk。
\tcblower
\silu 先假设“平衡不移动”，再与实际比较。\\
\jie 体积扩大为 2 倍的瞬间，W 浓度变为原来的 0.5 倍；最终为 0.65 倍 $>0.5$ 倍，说明平衡正向移动生成了更多 W。减压平衡向气体分子数增大的方向移动，故 $a+b<c+d$；X 的转化率增大。\\
\xianjing 若只看“0.65 倍 $<1$ 倍”就判断“逆移”，就犯了忽视体积变化的错误。
\end{liti}

\begin{yicuo}
下列事实\textbf{不能}用勒夏特列原理解释：使用催化剂加快合成氨；\ce{H2 + I2 <=> 2HI} 平衡体系加压颜色变深；工业合成氨选择 500 ℃左右（从平衡看低温才有利，选高温是为了速率和催化剂活性）。能解释的如：光照新制氯水 pH 减小（\ce{HClO} 分解，\ce{Cl2 + H2O <=> HCl + HClO} 正移）；浓氨水加 NaOH 固体制氨；增大压强有利于 \ce{SO2} 转化为 \ce{SO3}。
\end{yicuo}

\section{化学平衡图像}

\begin{fangfa}[图像题通用四步]
读图时依次看四处：一看面，纵、横坐标代表什么；二看线，曲线的变化趋势，是否为平衡线；三看点，起点、拐点、交点、最高（低）点、终点；四看辅助线，作等温线、等压线。分析时用两个基本原则。\textbf{①“先拐先平，数值大”}：先出现拐点（先达平衡）的曲线对应的温度高或压强大或使用了催化剂。\textbf{②“定一议二”}：图中有三个变量时，先固定一个，再讨论另外两个的关系。
\end{fangfa}

\subsection{速率—时间（$v$–$t$）图像}
\begin{center}
\begin{tikzpicture}[scale=0.55,font=\scriptsize]
\newcommand{\vtaxes}{\draw[-Stealth] (0,0)--(4.2,0) node[below]{$t$}; \draw[-Stealth] (0,0)--(0,3.2) node[left]{$v$}; \draw[thick] (0,1.3)--(1.2,1.3);}
\begin{scope}
\vtaxes \draw[thick,mainblue] (1.2,2.5) .. controls (2,1.9) .. (3.8,1.75); \draw[thick,warnred] (1.2,1.3) .. controls (2,1.6) .. (3.8,1.75); \draw[densely dotted] (1.2,1.3)--(1.2,2.5);
\node at (2.1,-0.6) {①增大反应物浓度};\node[mainblue] at (2.0,2.5){$\vz$};\node[warnred] at (2.2,1.25){$\vn$};
\end{scope}
\begin{scope}[xshift=5.2cm]
\vtaxes \draw[thick,mainblue] (1.2,2.0) .. controls (2,2.25) .. (3.8,2.35); \draw[thick,warnred] (1.2,2.7) .. controls (2,2.45) .. (3.8,2.35); \draw[densely dotted] (1.2,1.3)--(1.2,2.7);
\node[align=center] at (2.1,-0.8) {②升温\\（正反应放热）};\node[warnred] at (2.2,2.9){$\vn$};\node[mainblue] at (2.0,1.75){$\vz$};
\end{scope}
\begin{scope}[xshift=10.4cm]
\vtaxes \draw[thick,mainblue] (1.2,2.7) .. controls (2,2.45) .. (3.8,2.35); \draw[thick,warnred] (1.2,2.0) .. controls (2,2.25) .. (3.8,2.35); \draw[densely dotted] (1.2,1.3)--(1.2,2.7);
\node[align=center] at (2.1,-0.8) {③加压（正反应\\气体分子数减小）};\node[mainblue] at (2.2,2.9){$\vz$};\node[warnred] at (2.0,1.75){$\vn$};
\end{scope}
\begin{scope}[xshift=15.6cm]
\vtaxes \draw[thick,deepgreen] (1.2,2.4)--(3.8,2.4); \draw[densely dotted] (1.2,1.3)--(1.2,2.4);
\node[align=center] at (2.1,-0.8) {④催化剂\\（或等体积反应加压）};\node[deepgreen] at (2.5,2.75){$\vz'=\vn'$};
\end{scope}
\end{tikzpicture}
\end{center}
读图要点：①某一速率从原值开始连续变化，改变的是浓度；两条线都发生跳变（断点），改变的是温度或压强；两条线同步跳到同一值（平台），是使用了催化剂，或对等体积反应改变了压强。②条件改变后位于上方的那条速率线，决定平衡移动的方向。

\begin{liti}[多阶段 $v$-$t$ 图像]
合成氨反应 \ce{N2(g) + 3H2(g) <=> 2NH3(g)} $\Delta H<0$ 在某时段内的 $v$–$t$ 图如下（$t_1$、$t_3$、$t_4$ 时刻各改变一个条件）。(1) 判断各时刻改变的条件；(2) 哪段时间内 \ce{NH3} 的百分含量最高？
\begin{center}
\begin{tikzpicture}[xscale=0.95,yscale=0.8,font=\scriptsize]
\draw[-Stealth] (0,0)--(8.4,0) node[below]{$t$}; \draw[-Stealth] (0,0)--(0,3.6) node[left]{$v$};
\draw[thick] (0,1.2)--(1.2,1.2);
\draw[thick,warnred] (1.2,2.6) .. controls (1.8,2.2) .. (2.6,2.05); \draw[thick,mainblue] (1.2,1.8) .. controls (1.8,1.95) .. (2.6,2.05);
\draw[thick] (2.6,2.05)--(3.6,2.05);
\draw[thick,deepgreen] (3.6,2.9)--(5.2,2.9);
\draw[thick,mainblue] (5.2,1.4) .. controls (5.9,1.55) .. (6.8,1.62); \draw[thick,warnred] (5.2,1.9) .. controls (5.9,1.7) .. (6.8,1.62);
\draw[thick] (6.8,1.62)--(8.0,1.62);
\foreach \x/\l in {1.2/t_1,2.6/t_2,3.6/t_3,5.2/t_4,6.8/t_5}{\draw[densely dotted] (\x,0)--(\x,3.1); \node[below] at (\x,0) {$\l$};}
\node[warnred] at (1.65,2.8) {$\vn$}; \node[mainblue] at (1.65,1.55) {$\vz$};
\node[deepgreen] at (4.4,3.15) {$\vz=\vn$};
\node[warnred] at (5.65,2.05) {$\vn$}; \node[mainblue] at (5.65,1.25) {$\vz$};
\end{tikzpicture}
\end{center}
\tcblower
\jie (1) $t_1$：$\vz$、$\vn$ 都突然增大（断点），且 $\vn>\vz$，平衡逆移。若是加压，平衡应正移，所以是\textbf{升高温度}（放热反应升温逆移）。$t_3$：两速率同等程度增大（平台），平衡不移动，是\textbf{使用催化剂}（该反应气体分子数变化，不能是加压）。$t_4$：两速率都突然减小且 $\vn>\vz$，平衡逆移，是\textbf{减小压强}（若是降温，平衡应正移）。\\
(2) $t_1$ 后平衡逆移，$t_4$ 后又逆移，\ce{NH3} 含量逐次降低，所以 $0\sim t_1$ 段 \ce{NH3} 百分含量最高。\\
\xianjing 不要把“$t_3\sim t_4$ 速率最大”误认为“\ce{NH3} 含量最高”，速率大小与平衡组成是两回事。
\end{liti}

\subsection{含量（转化率）—时间—温度（压强）图像}
\begin{center}
\begin{tikzpicture}[scale=0.6,font=\scriptsize]
\begin{scope}
\draw[-Stealth] (0,0)--(5,0) node[below]{$t$}; \draw[-Stealth] (0,0)--(0,3.4) node[left]{C\%};
\draw[thick,warnred] (0,0) .. controls (0.6,1.6) .. (1.1,1.7)--(4.6,1.7) node[right]{$T_2$};
\draw[thick,mainblue] (0,0) .. controls (1.5,2.5) .. (2.4,2.7)--(4.6,2.7) node[right]{$T_1$};
\node[align=center] at (2.4,-0.9) {$T_2$ 先拐，$T_2>T_1$；高温 C\% 低\\$\Rightarrow$ 正反应放热};
\end{scope}
\begin{scope}[xshift=7.8cm]
\draw[-Stealth] (0,0)--(5,0) node[below]{$T$}; \draw[-Stealth] (0,0)--(0,3.4) node[left]{$\alpha$(A)};
\draw[thick,mainblue] (0.2,2.9) .. controls (2,2.6) .. (4.6,1.6) node[right]{$p_1$};
\draw[thick,warnred] (0.2,2.2) .. controls (2,1.9) .. (4.6,0.9) node[right]{$p_2$};
\draw[dashed] (2.5,0)--(2.5,3.0);
\node[align=center] at (2.4,-0.9) {升温 $\alpha$ 减小$\Rightarrow$放热；作等温线，\\若 $p_1>p_2$ 则正反应气体分子数减小};
\end{scope}
\end{tikzpicture}
\end{center}

\subsection{特殊点的分析}
\textbf{(1) 平衡曲线外的点：}平衡曲线上每一点都是平衡点。曲线外的点，看它要“走向”曲线需要正移还是逆移：若某点的生成物含量低于同条件下的平衡值，则反应还要正向进行，$\vz>\vn$。

\textbf{(2) 最高点（最低点）：}在“反应进行相同时间后测产物含量（或转化率）—温度”曲线中，若出现先升后降，则最高点之前多为“未达平衡”，升温加快速率，含量增大；最高点之后已达平衡，升温使平衡逆向移动（正反应放热），含量减小。也可能是催化剂在高温下活性降低所致，答题时要结合题给信息判断。

\begin{liti}[最高点图像]
在不同温度下，向相同的恒容密闭容器中充入等量 CO 和 \ce{H2}，发生 \ce{CO(g) + 2H2(g) <=> CH3OH(g)}，反应相同时间后测得 \ce{CH3OH} 的体积分数 $\varphi$ 与温度 $T$ 的关系如图（在 $T_3$ 处最高）。(1) $T_1$、$T_5$ 中，哪个点一定处于平衡状态？(2) 该反应 $\Delta H$\tk0。(3) $T_3$ 之前 $\varphi$ 随温度升高而增大的原因是\tk。
\begin{center}
\begin{tikzpicture}[scale=0.55,font=\scriptsize]
\draw[-Stealth] (0,0)--(6.5,0) node[below]{$T$}; \draw[-Stealth] (0,0)--(0,3.2) node[left]{$\varphi(\ce{CH3OH})$};
\draw[thick,mainblue] (0.4,0.4) .. controls (1.8,1.2) and (2.4,2.7) .. (3,2.7) .. controls (3.6,2.7) and (4.6,1.6) .. (6,1.0);
\foreach \x/\l in {0.9/1,2.0/2,3.0/3,4.4/4,5.6/5} {\draw[dashed] (\x,0)--(\x,0.12) node[below=2pt]{$T_\l$};}
\end{tikzpicture}
\end{center}
\tcblower
\jie (1) $T_5$。最高点之后的点（$T_4$、$T_5$）已达平衡；$T_1$ 时速率慢，相同时间内尚未达到平衡。(2) $<$：最高点之后升温 $\varphi$ 减小，平衡逆移，正反应放热。(3) $T_3$ 之前反应未达平衡，升高温度反应速率加快，相同时间内生成的 \ce{CH3OH} 增多。
\end{liti}

\section{化学反应的方向}

\subsection{自发过程与焓判据、熵判据}
在一定条件下，不借助外力就能自动进行的过程叫\textbf{自发过程}，如水往低处流、冰在常温下融化、墨水扩散。能自发进行到显著程度的反应叫\textbf{自发反应}。“自发”只说明有发生的趋势，不涉及快慢：\ce{H2} 与 \ce{O2} 常温下自发，但不点燃几乎不反应。

\textbf{焓判据：}放热反应体系能量降低，多数能自发。但有些吸热反应也自发，如 \ce{Ba(OH)2*8H2O} 与 \ce{NH4Cl} 的反应、\ce{NaHCO3} 与盐酸的反应。

\textbf{熵判据：}熵（$S$）是衡量体系混乱度的物理量。同一物质 $S(\mathrm{g})>S(\mathrm{l})>S(\mathrm{s})$；气体分子数增多的反应 $\Delta S>0$。熵增有利于自发，但有些熵减反应也自发，如 \ce{4Fe(OH)2(s) + O2(g) + 2H2O(l) = 4Fe(OH)3(s)}。

\begin{kanwu}[课件 2.3 第16页]
课件中“\ce{Fe(OH)2(s) + O2(g) + H2O(g) = Fe(OH)3(s)}”未配平且水的状态不当，应为 \ce{4Fe(OH)2(s) + O2(g) + 2H2O(l) = 4Fe(OH)3(s)}（这是氢氧化亚铁被氧化，而不是“制备氢氧化亚铁”）。
\end{kanwu}

\subsection{自由能判据}
\begin{gainian}[$\Delta G=\Delta H-T\Delta S$]
恒温恒压下：$\Delta G<0$，反应能自发进行；$\Delta G=0$，平衡状态；$\Delta G>0$，不能自发进行（逆反应自发）。
\end{gainian}

\begin{center}
\small
\begin{tabular}{cccl}
\toprule
$\Delta H$ & $\Delta S$ & $\Delta G=\Delta H-T\Delta S$ & 自发性 \\
\midrule
$<0$ & $>0$ & 恒 $<0$ & 任何温度都自发 \\
$>0$ & $<0$ & 恒 $>0$ & 任何温度都不自发 \\
$<0$ & $<0$ & 低温 $<0$ & 低温自发（如合成氨） \\
$>0$ & $>0$ & 高温 $<0$ & 高温自发（如 \ce{CaCO3} 分解） \\
\bottomrule
\end{tabular}
\end{center}
转折温度 $T=\dfrac{\Delta H}{\Delta S}$。计算时注意单位：$\Delta H$ 常用 kJ·mol$^{-1}$，$\Delta S$ 常用 J·mol$^{-1}$·K$^{-1}$，需统一。例如 \ce{CaCO3(s) = CaO(s) + CO2(g)}，$\Delta H=+178.5$ kJ·mol$^{-1}$，$\Delta S=+160.5$ J·mol$^{-1}$·K$^{-1}$，自发分解所需温度 $T>\dfrac{178.5\times10^3}{160.5}\approx1112$ K。

\begin{zhenti}{2025年广东卷第19题(2)\quad 来源：仓库所附真题文档}
298 K 下，反应 \ce{TiO2(s) + 2C(s) + 2Cl2(g) = TiCl4(g) + 2CO(g)} 的 $\Delta H<0$、$\Delta S>0$，则 298 K 下该反应\tk（填“能”或“不能”）自发进行。

\textbf{解析：}$\Delta H<0$、$\Delta S>0$，$\Delta G=\Delta H-T\Delta S$ 在任何温度下都小于 0，能自发进行。
\end{zhenti}

\begin{liti}[由自发性推断]
常温下，反应 \ce{A(s) + B(g) = 2X(g)} 不能自发进行。则该反应的 $\Delta H$\tk0。\ce{CH3OH(l) + NH3(g) = CH3NH2(g) + H2O(g)} 在某温度下自发向右进行，若 $|\Delta H|=17$ kJ·mol$^{-1}$，$|\Delta H-T\Delta S|=17$ kJ·mol$^{-1}$，则 $\Delta H$\tk0。
\tcblower
\jie 第一个反应气体分子数增大，$\Delta S>0$；若 $\Delta H\le0$，则 $\Delta G<0$ 一定自发，与事实矛盾，故 $\Delta H>0$。第二个反应自发，$\Delta G=-17$；该反应气体分子数增加，$\Delta S>0$，$-T\Delta S<0$。若 $\Delta H=-17$，则 $\Delta G=-17-T\Delta S<-17$，矛盾；故 $\Delta H=+17$ kJ·mol$^{-1}$，$\Delta H>0$。
\end{liti}

\section{化学反应的调控}

\subsection{工业合成氨条件的选择}
\ce{N2(g) + 3H2(g) <=> 2NH3(g)} $\Delta H=-92.4$ kJ·mol$^{-1}$。特点：可逆、放热、气体分子数减小（$\Delta S<0$，低温自发）。

\begin{center}
\small
\begin{tabularx}{\linewidth}{p{1.6cm}XXp{4.6cm}}
\toprule
条件 & 从速率看 & 从平衡看 & 实际选择及理由 \\
\midrule
压强 & 越大越快 & 越大越有利 & 10～30 MPa：压强越大，对设备材料和动力要求越高，成本增加 \\
温度 & 越高越快 & 越低越有利 & 400～500 ℃：兼顾速率；且铁触媒在 500 ℃ 左右活性最大 \\
催化剂 & 加快 & 无影响 & 铁触媒；原料气需净化，防止催化剂“中毒” \\
浓度 & 增大反应物浓度 & 增大反应物浓度、减小 \ce{NH3} 浓度 & 冷却使 \ce{NH3} 液化分离；未反应的 \ce{N2}、\ce{H2} 循环使用并及时补充；\ce{N2} 价廉，适当提高其比例以提高 \ce{H2} 的转化率 \\
\bottomrule
\end{tabularx}
\end{center}

\begin{center}
\small
\begin{tabular}{ccccccc}
\toprule
\multirow{2}{*}{温度/℃} & \multicolumn{6}{c}{平衡混合物中 \ce{NH3} 的体积分数/\%（$n(\ce{N2}):n(\ce{H2})=1:3$）} \\
\cmidrule(lr){2-7}
& 0.1 MPa & 10 MPa & 20 MPa & 30 MPa & 60 MPa & 100 MPa \\
\midrule
200 & 15.3 & 81.5 & 86.4 & 89.9 & 95.4 & 98.8 \\
300 & 2.20 & 52.0 & 64.2 & 71.0 & 84.2 & 92.6 \\
400 & 0.40 & 25.1 & 38.2 & 47.0 & 65.2 & 79.8 \\
500 & 0.10 & 10.6 & 19.1 & 26.4 & 42.2 & 57.5 \\
600 & 0.05 & 4.50 & 9.10 & 13.8 & 23.1 & 31.4 \\
\bottomrule
\end{tabular}
\end{center}
\noindent\small 数据分析：同一温度下压强越大、同一压强下温度越低，\ce{NH3} 的平衡含量越高，与勒夏特列原理的预测一致；但 200 ℃ 时速率太慢，实际只能折中选择 400～500 ℃，靠循环利用原料弥补单程转化率的不足。\normalsize

工艺流程中，原料气先经净化，防止催化剂中毒；再由压缩机加压；热交换器用反应放出的热预热原料气，节约能源；冷却使氨液化分离，促使平衡正移；未反应的原料气循环使用，提高原料利用率。

硫酸工业中 \ce{SO2} 的催化氧化可以用同样的思路分析：\ce{2SO2(g) + O2(g) <=> 2SO3(g)} $\Delta H=-196.6$ kJ·mol$^{-1}$。常压下 450 ℃ 时 \ce{SO2} 平衡转化率已达 97.5\%，加压收益很小而成本大，故采用常压；温度 400～500 ℃（兼顾速率与 \ce{V2O5} 催化剂活性）；通入过量空气提高 \ce{SO2} 转化率；尾气中的 \ce{SO2} 需回收处理。

\begin{fangfa}[“条件选择原因”的答题模板]
一般要从三个角度作答：①速率，该条件使反应速率较快；②平衡，该条件下平衡转化率（产率）较高或可接受；③实际因素，催化剂活性温度、设备承受能力、能耗与成本、安全与环保。例如“选择 500 ℃左右”：温度较高反应速率快，且铁触媒在 500 ℃左右活性最大；温度过高则平衡转化率降低、能耗增大。
\end{fangfa}

\section{平衡计算的常见题型}

\begin{fangfa}[五类常见平衡计算]
①求转化率、平衡浓度：三段式，转化量成比例。\\
②已知 $K$ 求组成：设未知数代入 $K$ 表达式，常解一元二次方程；若 $K$ 很大或很小，可先假设“几乎完全反应”或“几乎不反应”再近似。\\
③压强、密度、平均摩尔质量的变化：恒温恒容 $\dfrac{p_\text{后}}{p_\text{前}}=\dfrac{n_\text{后}}{n_\text{前}}$；恒温恒压 $\dfrac{V_\text{后}}{V_\text{前}}=\dfrac{n_\text{后}}{n_\text{前}}$；$\bar M=\dfrac{m_\text{总}}{n_\text{总}}$。\\
④平衡后再加料：先判断是恒容还是恒压，用“虚拟容器法”转化为等效平衡或加压问题。\\
⑤范围（极值）问题：可逆反应不能进行到底，各物质的量介于“完全正向”与“完全逆向”两个极端之间。
\end{fangfa}

\begin{liti}[极值法判断范围]
在密闭容器中进行 \ce{X2(g) + Y2(g) <=> 2Z(g)}，已知 \ce{X2}、\ce{Y2}、Z 的起始浓度分别为 0.1、0.3、0.2 mol·L$^{-1}$，一定条件下达到平衡。下列浓度可能存在的是
\xx{Z 为 0.4 mol·L$^{-1}$}{\ce{X2} 为 0.2 mol·L$^{-1}$}{\ce{Y2} 为 0.4 mol·L$^{-1}$}{Z 为 0.3 mol·L$^{-1}$}
\tcblower
\jie 假设完全正向：\ce{X2} 0，\ce{Y2} 0.2，Z 0.4；假设完全逆向：\ce{X2} 0.2，\ce{Y2} 0.4，Z 0。可逆反应不能进行到底，所以 $0<c(\ce{X2})<0.2$，$0.2<c(\ce{Y2})<0.4$，$0<c(\mathrm{Z})<0.4$。A、B、C 都取到了端点值，不可能；D 可能。答案：D。
\end{liti}

\begin{liti}[由压强变化求转化率]
恒温恒容下，向密闭容器中充入 2 mol \ce{SO2} 和 1 mol \ce{O2}，发生 \ce{2SO2(g) + O2(g) <=> 2SO3(g)}，平衡时容器内压强为起始时的 $\frac56$。求 \ce{SO2} 的平衡转化率及平衡时 \ce{SO3} 的体积分数。
\tcblower
\jie 设转化 \ce{O2} $x$ mol，则气体总物质的量减少 $x$ mol：$\dfrac{3-x}{3}=\dfrac56$，$x=0.5$。转化 \ce{SO2} 1 mol，$\alpha(\ce{SO2})=50\%$；平衡时 \ce{SO3} 1 mol，总 2.5 mol，体积分数 40\%。
\end{liti}

\section{多重平衡与选择性}

实际体系中往往同时存在几个反应，它们共用某些物质（\textbf{竞争反应}或\textbf{连续反应}）。处理多重平衡时注意以下几点：
\begin{enumerate}
\item 每个反应各设一个变量（转化量 $x$、$y$），列出各物质的平衡量；共用物质的量要把各反应的贡献相加。
\item 用原子守恒代替繁琐的三段式往往更快：如由 C 原子守恒得 $n(\ce{CO2})+n(\ce{CO})+n(\ce{CH3OH})=n_0(\ce{CO2})$。
\item 每个反应的 $K$ 都要用同一平衡体系中各物质的平衡浓度（分压）计算，共用物质的浓度只有一个。
\item 分析温度对转化率的影响时，要看哪个反应占主导：放热反应在低温段占主导，吸热反应在高温段占主导。
\end{enumerate}

\begin{liti}[多重平衡的计算]
\ce{CO2} 加氢体系：反应Ⅰ \ce{CO2(g) + 3H2(g) <=> CH3OH(g) + H2O(g)} $\Delta H_1=-49.5$ kJ·mol$^{-1}$；反应Ⅱ \ce{CO2(g) + H2(g) <=> CO(g) + H2O(g)} $\Delta H_2=+41.2$ kJ·mol$^{-1}$。恒压 $p$ 下充入 1 mol \ce{CO2} 和 3 mol \ce{H2}，平衡时 \ce{CO2} 转化率为 30\%，\ce{CH3OH} 的选择性为 80\%。(1) 求平衡时各物质的物质的量；(2) 求反应Ⅱ的 $K_p$；(3) 实验发现：温度升高，\ce{CO2} 的平衡转化率先减小后增大，解释原因。
\tcblower
\jie (1) 转化的 \ce{CO2} 为 0.3 mol，其中 $0.3\times80\%=0.24$ mol 生成 \ce{CH3OH}，0.06 mol 生成 CO。\ce{H2} 消耗 $3\times0.24+0.06=0.78$ mol，剩 2.22 mol；\ce{H2O} 由两反应共生成 $0.24+0.06=0.30$ mol。平衡时：\ce{CO2} 0.70，\ce{H2} 2.22，\ce{CH3OH} 0.24，\ce{H2O} 0.30，CO 0.06，$n_\text{总}=3.52$ mol。\\
(2) 反应Ⅱ前后气体分子数相等，$K_p=\dfrac{n(\ce{CO})n(\ce{H2O})}{n(\ce{CO2})n(\ce{H2})}=\dfrac{0.06\times0.30}{0.70\times2.22}\approx0.0116$。反应Ⅰ前后气体分子数不等，需代入分压（$p$ 为平衡总压）：\[K_p(\text{Ⅰ})=\dfrac{\dfrac{0.24}{3.52}p\cdot\dfrac{0.30}{3.52}p}{\dfrac{0.70}{3.52}p\cdot\left(\dfrac{2.22}{3.52}p\right)^3}\approx\dfrac{0.116}{p^2}\]
(3) 反应Ⅰ放热、反应Ⅱ吸热。低温段以反应Ⅰ为主，升温反应Ⅰ逆移，\ce{CO2} 转化率减小；高温段以反应Ⅱ为主，升温反应Ⅱ正移，\ce{CO2} 转化率增大。同时 \ce{CH3OH} 的选择性随温度升高持续下降。
\end{liti}

\begin{zhenti}{2025年广东卷第19题(3)②③\quad 来源：仓库所附真题文档}
将 5.0 mol \ce{TiCl4} 与 10.0 mol Ti 放入容积为 $V_0$ L 的恒容密闭容器中，体系存在：(a)\ce{Ti(s) + TiCl4(g) <=> 2TiCl2(s)} $\Delta H_1=-267.8$ kJ·mol$^{-1}$；(b)\ce{TiCl2(s) + TiCl4(g) <=> 2TiCl3(g)} $\Delta H=+200.1$ kJ·mol$^{-1}$；(c)\ce{Ti(s) + 3TiCl4(g) <=> 4TiCl3(g)} $\Delta H=+132.4$ kJ·mol$^{-1}$；(d)\ce{TiCl3(g) <=> TiCl3(s)} $\Delta H_2$（凝华，放热）。不同温度下平衡体系组成如图（1073～1273 K）：Ti(s) 的量随温度升高而增大，$T_0$ 时为 8.5 mol；曲线Ⅱ随温度升高而增大，$T_0$ 时为 6.0 mol；曲线Ⅰ随温度升高由约 9 mol 减小至接近 0；\ce{TiCl3(s)}、\ce{TiCl4(g)} 的量均很少。②曲线Ⅰ对应的物质为\tk。③温度 $T_0$ K 下（此时 \ce{TiCl2(s)}、\ce{TiCl3(s)} 均为 0），$n[\ce{TiCl4(g)}]=$\tk mol，反应(c)的平衡常数 $K=$\tk（列出算式）。

\textbf{解析：}②(a)放热、(d)放热，升温逆移；(b)(c)吸热，升温正移。所以 \ce{TiCl3(g)} 随温度升高增多（曲线Ⅱ），\ce{TiCl2(s)} 随温度升高减少（曲线Ⅰ）。③由 Ti 原子守恒：$n(\ce{TiCl4})=10.0+5.0-8.5-6.0=0.5$ mol（可用 Cl 守恒验证：$4\times0.5+3\times6.0=20=4\times5.0$）。Ti 为固体不写入，$K=\dfrac{c^4(\ce{TiCl3})}{c^3(\ce{TiCl4})}=\dfrac{(6.0/V_0)^4}{(0.5/V_0)^3}$。
\end{zhenti}

% ---------------- 第3章 章末 ----------------
\zhangmo{本章小结}
\subsection*{一、知识总结}
\begin{center}
\small
\begin{tabularx}{\linewidth}{p{2.4cm}X}
\toprule
平衡状态 & 逆、等、动、定、变；判断：正逆相等（一正一逆、成比例）或变量不变 \\
平衡常数 & 只与温度有关；固体纯液体不写；方程式相加 $K$ 相乘；$Q$ 与 $K$ 比较定方向；$K_p$ 用分压 \\
平衡移动 & 浓度、压强（看气体分子数）、温度（唯一改变 $K$）、催化剂（不移动）；勒夏特列原理“减弱不抵消” \\
图像 & 四看；先拐先平数值大；定一议二；曲线外的点看趋向；最高点前未平衡、后平衡 \\
方向 & $\Delta G=\Delta H-T\Delta S<0$ 自发；四种组合；转折温度 $\Delta H/\Delta S$ \\
调控 & 合成氨 10～30 MPa、400～500 ℃、铁触媒、液化分离与循环；兼顾速率、平衡、成本 \\
多重平衡 & 每个反应设变量；原子守恒；同一体系共用浓度；选择性$=$生成目标产物所耗反应物/反应物总转化量 \\
\bottomrule
\end{tabularx}
\end{center}

\subsection*{二、方法提升}
\begin{enumerate}
\item \textbf{三段式是平衡计算的“万能钥匙”}：起始、转化、平衡三行，转化量与计量数成比例；有 $K_p$ 时加一行“分压”。
\item \textbf{判断移动方向先用 $Q$ 与 $K$}：浓度、压强变化看 $Q$ 怎么变，温度变化看 $K$ 怎么变。
\item \textbf{“假设不移动再比较”}：处理“浓度变为原来的几倍”类问题，先算假设不移动时的值，再与实际值比较。
\item \textbf{多重平衡温度曲线}：先判断各反应的吸放热，再分温度段讨论谁占主导。
\end{enumerate}

\subsection*{三、易错点归纳}
\begin{enumerate}
\item 把“浓度相等”“浓度比等于计量数比”当成平衡标志；忘记检查一个量是否为“变量”。
\item $K$ 表达式中写入固体、纯液体；用起始浓度或某时刻浓度代替平衡浓度计算 $K$。
\item 误认为增大压强平衡一定移动；恒容充惰性气体判为“加压，平衡移动”。
\item 误认为平衡正移，反应物转化率一定增大（增大某反应物浓度时其自身转化率减小）。
\item 误认为“自发反应一定能迅速发生”；计算 $T=\Delta H/\Delta S$ 时忘统一 kJ 与 J。
\item 选择合成氨温度时答“温度高有利于平衡正移”（错，低温才有利于平衡）。
\end{enumerate}

\zhangmo{分层练习}
\zu{A组\quad 基础巩固}
\begin{xiti}
\item 一定温度下的恒容容器中发生 \ce{2A(g) + B(s) <=> C(g) + D(g)}，下列描述能表明反应已达平衡的是\xz\\
①混合气体的压强不变\quad ②混合气体的密度不变\quad ③C(g) 的浓度不变\quad ④A、C、D 的浓度之比为 2:1:1\quad ⑤单位时间内生成 $n$ mol C，同时生成 $n$ mol D\quad ⑥单位时间内生成 $n$ mol D，同时生成 $2n$ mol A
\xx{①②⑤}{②③⑤}{②③⑥}{①③⑥}

\item 下列关于化学平衡常数的说法正确的是\xz
\xxa{化学平衡常数表示反应进行的程度，与温度无关}{两种物质反应，不管怎样书写化学方程式，平衡常数均不变}{化学平衡常数等于某一时刻生成物浓度幂之积与反应物浓度幂之积的比值}{温度一定时，对于给定的反应，正、逆反应的平衡常数互为倒数}

\item 某温度下 \ce{N2(g) + 3H2(g) <=> 2NH3(g)} 的平衡常数为 $K$。同温下，\ce{1/2N2(g) + 3/2H2(g) <=> NH3(g)} 的平衡常数为\tk，\ce{2NH3(g) <=> N2(g) + 3H2(g)} 的平衡常数为\tk。

\item 下列事实\textbf{不能}用勒夏特列原理解释的是\xz
\xxa{光照新制的氯水时，溶液的 pH 逐渐减小}{加催化剂，使 \ce{N2} 和 \ce{H2} 在一定条件下转化为 \ce{NH3}}{可用浓氨水和氢氧化钠固体快速制氨}{增大压强，有利于 \ce{SO2} 与 \ce{O2} 反应生成 \ce{SO3}}

\item 对于 \ce{2SO2(g) + O2(g) <=> 2SO3(g)} $\Delta H<0$，在恒温恒容条件下达到平衡后，下列措施能提高 \ce{SO2} 平衡转化率的是\xz\\
①升高温度\quad ②再充入 \ce{O2}\quad ③充入 He\quad ④再按原比例充入 \ce{SO2} 和 \ce{O2}\quad ⑤使用催化剂\quad ⑥移走部分 \ce{SO3}
\xx{①②⑥}{②④⑥}{②③④}{④⑤⑥}

\item 下列反应在任何温度下均能自发进行的是\xz
\xxa{\ce{2N2(g) + O2(g) = 2N2O(g)} \quad $\Delta H=+163$ kJ·mol$^{-1}$}{\ce{Ag(s) + 1/2Cl2(g) = AgCl(s)} \quad $\Delta H=-127$ kJ·mol$^{-1}$}{\ce{HgO(s) = Hg(l) + 1/2O2(g)} \quad $\Delta H=+91$ kJ·mol$^{-1}$}{\ce{H2O2(l) = 1/2O2(g) + H2O(l)} \quad $\Delta H=-98$ kJ·mol$^{-1}$}

\item 在 \ce{2NO2(g) <=> N2O4(g)} $\Delta H<0$ 的平衡体系中，下列操作能使混合气体颜色变深、且平衡逆向移动的是\xz
\xxa{恒温时压缩容器体积}{恒容时升高温度}{恒容时充入 Ar}{恒压时充入 Ar}
\end{xiti}

\zu{B组\quad 高考典型}
\begin{xiti}
\item 某温度下，在 2 L 密闭容器中充入气体 A、B，其浓度分别为 2 mol·L$^{-1}$、1 mol·L$^{-1}$，发生 \ce{3A(g) + 2B(g) <=> 4C(\text{?}) + 2D(\text{?})}。达到平衡时生成 1.6 mol C，且反应前后的压强之比为 5:4（相同温度）。(1) 判断 C、D 的状态，写出平衡常数表达式；(2) 求 B 的平衡转化率和该温度下的 $K$（保留两位有效数字）；(3) 增加 C 的量，A、B 的转化率如何变化？

\item \laiyuan{原创}\ 对于 \ce{$m$A(g) + $n$B(g) <=> $p$C(g)}，在不同条件下 C 的体积分数随时间变化如下：曲线 a（$T_2$、$p_2$）最先出现拐点，平台高度 0.30；曲线 b（$T_1$、$p_2$）其次出现拐点，平台高度 0.50；曲线 c（$T_1$、$p_1$）最后出现拐点，平台高度 0.35。(1) 比较 $T_1$、$T_2$ 和 $p_1$、$p_2$ 的大小；(2) 判断 $\Delta H$ 的正负及 $m+n$ 与 $p$ 的关系；(3) 在恒压下，\ce{2SO2(g) + O2(g) <=> 2SO3(g)} 的 \ce{SO2} 平衡转化率随温度升高而降低，投料比 $m=n(\ce{SO2})/n(\ce{O2})$ 分别为 $m_1$、$m_2$、$m_3$ 的三条曲线自上而下排列，比较 $m_1$、$m_2$、$m_3$ 的大小。

\item 硫酸工业中 \ce{2SO2(g) + O2(g) <=> 2SO3(g)} $\Delta H=-196.6$ kJ·mol$^{-1}$。不同条件下 \ce{SO2} 的平衡转化率：450 ℃ 时 0.1 MPa 为 97.5\%，10 MPa 为 99.7\%；550 ℃ 时 0.1 MPa 为 85.6\%，10 MPa 为 98.3\%。(1) 实际生产为何选择常压？(2) 为何选择 400～500 ℃ 而不是更低温度？(3) 为何通入过量空气？(4) 尾气为何要处理？

\item \laiyuan{2024年广东卷第17题(4)，来源：仓库所附真题文档}\ 兴趣小组探究反应 \ce{CaSO4(s) + CO3^2-(aq) <=> CaCO3(s) + SO4^2-(aq)}（反应Ⅰ）平衡的建立：将 1.00 g 硫酸钙（$M=136$ g·mol$^{-1}$）加入 100.0 mL 0.100 mol·L$^{-1}$ \ce{Na2CO3} 溶液中，在 25 ℃ 和搅拌条件下，用 pH 计测得体系 pH 随时间先减小，$t_2$ min 后不变。甲同学认为反应Ⅰ已达平衡；乙同学认为证据不足，提出：假设1 硫酸钙固体已完全消耗；假设2 硫酸钙有剩余，但被碳酸钙沉淀包裹。(1) 乙设计实验：将反应混合物固液分离，取少量沉淀置于试管中，滴加\tk，现象为\tk，沉淀完全溶解；继续向试管中滴加\tk，无白色沉淀生成。由此得出假设1成立。(2) 从数据角度说明假设1为何成立。(3) 写出优化的实验方案，并给出反应Ⅰ平衡已建立的判断依据。

\item \laiyuan{原创}\ 在恒定总压 20 MPa、一定温度下，向反应器中充入 1 mol \ce{N2} 和 3 mol \ce{H2}，发生 \ce{N2(g) + 3H2(g) <=> 2NH3(g)}，平衡时 \ce{NH3} 的物质的量分数为 25\%。(1) 求 \ce{H2} 的平衡转化率。(2) 求该温度下的 $K_p$（保留三位有效数字，单位 MPa$^{-2}$）。(3) 若改为在恒温恒容条件下进行（起始压强仍为 20 MPa），\ce{NH3} 的平衡物质的量分数将\tk（填“大于”“小于”或“等于”）25\%。

\item \laiyuan{原创}\ \ce{CO2(g) + 3H2(g) <=> CH3OH(g) + H2O(g)} $\Delta H<0$。在压强一定时，按不同投料比 $m=\dfrac{n(\ce{H2})}{n(\ce{CO2})}$（$m_1$、$m_2$、$m_3$）进行实验，\ce{CO2} 的平衡转化率随温度变化的三条曲线均随温度升高而下降，且同一温度下 $m_1$ 对应的转化率最高、$m_3$ 最低。(1) 比较 $m_1$、$m_2$、$m_3$ 的大小。(2) 解释曲线随温度升高而下降的原因。(3) 若同一投料比下，曲线上 a 点（$T_a$）与 b 点（$T_b$，$T_b>T_a$）的平衡常数分别为 $K_a$、$K_b$，比较大小。
\end{xiti}

\zu{C组\quad 冲刺高分}
\begin{xiti}
\item \laiyuan{原创}\ 乙苯脱氢制苯乙烯：\ce{C6H5CH2CH3(g) <=> C6H5CH=CH2(g) + H2(g)} $\Delta H>0$。某温度、总压恒为 100 kPa 时，以纯乙苯为原料，乙苯的平衡转化率为 40\%。(1) 求该温度下的 $K_p$。(2) 保持温度和总压不变，改为通入 $n(\ce{H2O}):n(\text{乙苯})=9:1$ 的混合气（水蒸气不参与反应），求乙苯的平衡转化率（保留两位有效数字）。(3) 解释工业上掺入大量水蒸气的原因（至少答两点）。

\item \laiyuan{2025年广东卷第19题(4)，来源：仓库所附真题文档}\ 钛基催化剂可催化储氢物质肼（\ce{N2H4}）的分解：(e)\ce{N2H4 = N2 + 2H2}；(f)\ce{3N2H4 = N2 + 4NH3}。以肼的水溶液为原料（含 \ce{N2H4} 的物质的量为 $n_0$）进行实验，得到 $n_1/n_0$、$n_2/n_0$ 随时间 $t$ 的变化曲线，其中 $n_1$ 为 \ce{H2} 与 \ce{N2} 的物质的量之和，$n_2$ 为剩余 \ce{N2H4} 的物质的量；$t_0$ min 时 $n_1/n_0=y_1$，$n_2/n_0=y_2$。设 $n_\text{e}$ 为 $0\sim t$ 时间段内反应(e)消耗 \ce{N2H4} 的物质的量，该时间段内催化剂的选择性用 $\dfrac{n_\text{e}}{n_0-n_2}\times100\%$ 表示。①$0\sim t_0$ min 内，\ce{N2H4} 的转化率为\tk（用含 $y_2$ 的代数式表示）。②$0\sim t_0$ min 内，催化剂的选择性为\tk（用含 $y_1$ 与 $y_2$ 的代数式表示，写出推导过程）。

\item \laiyuan{2026年广东卷第19题(5)，来源：仓库所附真题文档（网络收集版）}\ 某低共熔溶剂 D(l) 可用于吸收 HCl(g)，吸收过程可表示为 \ce{HCl(g) <=> HCl(sol)}，\ce{$x$HCl(sol) + D(l) <=> D*$x$HCl(sol)}。在一定 $T$ 和 HCl(g) 压强 $p$ 下，吸收至饱和时，每摩尔 D 吸收 HCl(g) 的量为 $\alpha+\beta$，其中 $\alpha=\dfrac{n(\ce{HCl})}{n_0(\mathrm{D})}$，$\beta=\dfrac{x\cdot n(\mathrm{DH})}{n_0(\mathrm{D})}$（$x$ 为仅与温度有关的系数；$n_0(\mathrm{D})$ 为 D(l) 的初始物质的量；$n(\ce{HCl})$ 为平衡时 HCl(sol) 的物质的量；$n(\mathrm{DH})$ 为平衡时 \ce{D*$x$HCl(sol)} 的物质的量）。由图读得：303 K、$p=100$ kPa 时 $\alpha=0.15$、$\beta=0.86$；350 K 时 $\alpha$ 始终为 0.00，$p=40$ kPa 时 $\beta=0.19$，$p=100$ kPa 时 $\beta=0.25$。
①1 mol D 在 303 K、$p=100$ kPa 下吸收 HCl(g) 至饱和，再升温至 350 K、降压至 $p=40$ kPa，达平衡后可释放\tk mol HCl(g)。
②形成 \ce{D*$x$HCl(sol)} 的总过程为 \ce{$x$HCl(g) + D(l) <=> D*$x$HCl(sol)}，其平衡常数 $K=\dfrac{n(\mathrm{DH})}{n(\mathrm{D})\cdot(p/100\,\mathrm{kPa})^x}$。已知 350 K 时 $K=1.0$，则在 350 K、$p=64$ kPa 下，1 mol D 最多可吸收\tk mol HCl(g)（写出计算过程，结果保留两位有效数字）。

\item \laiyuan{原创}\ 一定温度下，在恒容密闭容器中充入 \ce{NO2}，起始压强为 100 kPa，发生 \ce{2NO2(g) <=> N2O4(g)}，该反应可视为基元反应：$\vz=\kz p^2(\ce{NO2})$，$\vn=\kn p(\ce{N2O4})$。达到平衡时总压为 70 kPa。(1) 求平衡时 \ce{NO2}、\ce{N2O4} 的分压及 $K_p$。(2) 用 $\kz$、$\kn$ 表示 $K_p$。(3) 当 \ce{NO2} 的分压为 60 kPa 时，$\dfrac{\vz}{\vn}=$\tk。(4) 升高温度，平衡体系中 \ce{NO2} 的体积分数增大，判断该反应的 $\Delta H$ 正负。

\item \laiyuan{改编自选必1课件“演示文稿1”中的练习，数据为本书设定}\ 羰基法提纯镍：\ce{Ni(s) + 4CO(g) <=> Ni(CO)4(g)}。在不同总压下，平衡时 \ce{Ni(CO)4} 的体积分数 $x$ 随温度升高而减小，且同一温度下总压越大 $x$ 越大。(1) 判断该反应的 $\Delta H$\tk0，从热力学角度考虑，\tk{}有利于 \ce{Ni(CO)4} 的生成（写出两点）。(2) 设起始只充入 CO（Ni 足量），若总压为 1.0 MPa、100 ℃ 时 $x=0.90$，求 CO 的平衡转化率（保留三位有效数字）和该温度下的 $K_p$（单位 MPa$^{-3}$）。(3) 工业上先在约 50 ℃ 使粗镍与 CO 反应生成 \ce{Ni(CO)4}，再将气体送到约 230 ℃ 的分解器得到纯镍。结合 (1) 说明其原理。
\end{xiti}

\zhangmo{练习解析}
\begin{jiexilist}
\item \ans{C}B 为固体，气体分子数 $2\to2$ 不变，恒容时压强始终不变，①不能作标志；气体质量随反应变化，密度是变量，②能；③浓度不变，能；④浓度比不能说明不变；⑤都表示正反应；⑥生成 D 为正、生成 A 为逆，$n:2n=1:2$ 与计量数比相符，能。选②③⑥。
\item \ans{D}$K$ 与温度有关，A 错；方程式写法不同 $K$ 不同，B 错；$K$ 用平衡浓度，C 错；正逆反应 $K$ 互为倒数，D 对。
\item \ans{$K^{1/2}$（即 $\sqrt K$）；$1/K$}计量数变为 $\frac12$，$K$ 开平方；逆反应取倒数。
\item \ans{B}催化剂不使平衡移动，不能用勒夏特列原理解释。A：\ce{HClO} 见光分解，\ce{Cl2 + H2O <=> H+ + Cl- + HClO} 正移，$c(\ce{H+})$ 增大；C：NaOH 溶解放热且 $c(\ce{OH-})$ 增大，使 \ce{NH3 + H2O <=> NH3*H2O <=> NH4+ + OH-} 逆移；D：加压平衡向气体分子数减小的方向移动。
\item \ans{B}①升温平衡逆移，转化率降低；②增大 \ce{O2} 浓度，\ce{SO2} 转化率提高；③恒容充 He，平衡不移动；④按比例加料相当于加压，平衡正移，转化率提高；⑤催化剂不影响平衡；⑥移走生成物，平衡正移。选②④⑥。
\item \ans{D}任何温度自发需 $\Delta H<0$ 且 $\Delta S>0$。A：吸热、气体分子数减少（熵减），任何温度不自发；B：放热、熵减，低温自发；C：吸热、熵增，高温自发；D：放热、生成气体熵增，任何温度自发。
\item \ans{B}A：压缩体积，\ce{NO2} 浓度瞬间增大，颜色变深，但平衡正向（气体分子数减小的方向）移动；B：升温，平衡向吸热的逆反应方向移动，\ce{NO2} 增多，颜色变深，正确；C：恒容充 Ar，各物质浓度不变，平衡不移动，颜色不变；D：恒压充 Ar，体积增大，相当于减压，平衡逆向移动，但 \ce{NO2} 浓度减小，颜色变浅。
\item (1) 起始 A 4 mol、B 2 mol（共 6 mol 气体）；转化 A 1.2 mol、B 0.8 mol、C 1.6 mol、D 0.8 mol。平衡气体总量 $6\times\frac45=4.8$ mol。若 C、D 均为气体，平衡时 $6-2.0+2.4=6.4$ mol，不符；若只有 D 为气体，$6-2.0+0.8=4.8$ mol，符合。故 C 为非气态（固体或液体），D 为气体，$K=\dfrac{c^2(\mathrm{D})}{c^3(\mathrm{A})\cdot c^2(\mathrm{B})}$。(2) $\alpha(\mathrm{B})=0.8/2=40\%$；平衡浓度 $c(\mathrm{A})=1.4$、$c(\mathrm{B})=0.6$、$c(\mathrm{D})=0.4$ mol·L$^{-1}$，$K=\dfrac{0.4^2}{1.4^3\times0.6^2}\approx0.16$。(3) C 非气态，增加 C 的量平衡不移动，A、B 转化率不变。
\item (1) 曲线 a、b 压强相同，a 先拐，$T_2>T_1$；曲线 b、c 温度相同，b 先拐，$p_2>p_1$。(2) 同为 $p_2$，高温（a）C\% 低，升温平衡逆移，$\Delta H<0$；同为 $T_1$，高压（b）C\% 高，加压平衡正移，$m+n>p$。(3) 投料比 $m$ 越小，\ce{O2} 相对越多，\ce{SO2} 转化率越高，最上方曲线 $m$ 最小：$m_1<m_2<m_3$。
\item (1) 常压下 \ce{SO2} 平衡转化率已很高（450 ℃ 时 97.5\%），加压提高有限，而对设备要求和能耗大幅增加，不经济。(2) 温度过低反应速率太慢，且催化剂（\ce{V2O5}）在 400～500 ℃ 活性最高；温度过高平衡转化率下降。(3) 增大 \ce{O2} 浓度，使平衡正移，提高 \ce{SO2} 的转化率（空气廉价易得）。(4) \ce{SO2} 有毒、会形成酸雨污染环境，回收还可提高原料利用率。
\item (1) 过量稀盐酸；有气体产生；\ce{BaCl2} 溶液。（\ce{CaCO3} 溶于盐酸放出 \ce{CO2}；若沉淀中有 \ce{CaSO4}，溶解后加 \ce{BaCl2} 会生成 \ce{BaSO4} 白色沉淀，无沉淀说明无 \ce{CaSO4}。）(2) $n(\ce{CaSO4})=1.00/136\approx7.4\times10^{-3}$ mol，$n(\ce{CO3^2-})=0.100\times0.100=1.00\times10^{-2}$ mol，\ce{CaSO4} 不足；反应Ⅰ的平衡常数很大，\ce{CaSO4} 可完全转化。pH 不变只是因为一种反应物耗尽，并非平衡。(3) 要保证 \ce{CaSO4} 有剩余：按完全反应计算至少需 $0.0100\times136=1.36$ g，故将不少于 1.36 g 硫酸钙加入 100.0 mL 0.100 mol·L$^{-1}$ \ce{Na2CO3} 溶液中，在 25 ℃、搅拌条件下用 pH 计测 pH；当 pH 不变时，加入 \ce{Na2SO4} 固体，若 pH 增大（\ce{SO4^2-} 浓度增大使反应Ⅰ逆移，$c(\ce{CO3^2-})$ 增大，碱性增强），说明反应Ⅰ的平衡已建立（体现了“平衡可以被打破并移动”这一特征）。
\item (1) 设 \ce{N2} 转化 $x$ mol：平衡时 \ce{N2} $(1-x)$、\ce{H2} $(3-3x)$、\ce{NH3} $2x$，总 $(4-2x)$ mol。$\dfrac{2x}{4-2x}=0.25$，$x=0.4$；$\alpha(\ce{H2})=\dfrac{1.2}{3}=40\%$。(2) $n_\text{总}=3.2$ mol，物质的量分数：\ce{N2} 0.1875、\ce{H2} 0.5625、\ce{NH3} 0.25。$K_p=\dfrac{(0.25\times20)^2}{(0.1875\times20)\times(0.5625\times20)^3}\approx4.68\times10^{-3}$ MPa$^{-2}$。(3) 小于。恒容时随反应进行气体分子数减少、压强降低，相当于在比恒压时更低的压强下建立平衡，平衡转化率较低。
\item (1) $m_1>m_2>m_3$：\ce{H2} 相对越多，\ce{CO2} 的转化率越高。(2) 正反应放热，升温平衡逆向移动，\ce{CO2} 平衡转化率下降。(3) $K_a>K_b$（放热反应，温度越高 $K$ 越小；$K$ 只与温度有关，与投料比无关）。
\item (1) 设乙苯 1 mol，平衡时乙苯 0.6 mol、苯乙烯 0.4 mol、\ce{H2} 0.4 mol，$n_\text{总}=1.4$ mol。$K_p=\dfrac{(\frac{0.4}{1.4}p)^2}{\frac{0.6}{1.4}p}=\dfrac{0.16}{0.84}\times100\approx19$ kPa。(2) 设转化率为 $\alpha$，$n_\text{总}=10+\alpha$：$K_p=\dfrac{\alpha^2p}{(1-\alpha)(10+\alpha)}=19.05$，整理得 $119.05\alpha^2+171.45\alpha-190.5=0$，$\alpha\approx0.74$，即约 74\%（远高于 40\%）。(3) ①恒压下水蒸气稀释，降低各反应物、产物的分压，相当于减压，平衡向气体分子数增大的正方向移动，提高乙苯转化率；②水蒸气热容大，可为吸热反应供热、维持反应温度；③水蒸气可与催化剂表面积炭反应（\ce{C + H2O -> CO + H2}），减少积炭，延长催化剂寿命。
\item ①$(1-y_2)\times100\%$。起始 $n_2/n_0=1$，$t_0$ 时为 $y_2$，转化了 $n_0(1-y_2)$。②设反应(e)消耗 \ce{N2H4} $n_\text{e}$，反应(f)消耗 $n_\text{f}$：$n_\text{e}+n_\text{f}=n_0(1-y_2)$；(e)生成 $n_\text{e}$ mol \ce{N2} 和 $2n_\text{e}$ mol \ce{H2}，(f)生成 $\frac13n_\text{f}$ mol \ce{N2}：$3n_\text{e}+\frac13n_\text{f}=y_1n_0$。将第一式乘 $\frac13$ 后与第二式相减：$\frac83n_\text{e}=n_0\left(y_1-\frac{1-y_2}{3}\right)$，$n_\text{e}=\dfrac{(3y_1+y_2-1)n_0}{8}$。选择性 $=\dfrac{n_\text{e}}{n_0(1-y_2)}\times100\%=\dfrac{3y_1+y_2-1}{8-8y_2}\times100\%$。
\item ①0.82。303 K、100 kPa 饱和时每摩尔 D 吸收 $\alpha+\beta=0.15+0.86=1.01$ mol；350 K、40 kPa 平衡时为 $0+0.19=0.19$ mol；释放 $1.01-0.19=0.82$ mol。②0.22。350 K 时 $\alpha=0$，HCl 只以 \ce{D*$x$HCl} 形式被吸收。$p=100$ kPa 时 $\beta=0.25$：$n(\mathrm{DH})=\dfrac{0.25}{x}$ mol，$n(\mathrm{D})=1-\dfrac{0.25}{x}$，$K=\dfrac{0.25/x}{(1-0.25/x)\times1^x}=1.0$，解得 $\dfrac{0.25}{x}=0.5$，$x=0.5$。$p=64$ kPa 时：$\dfrac{n(\mathrm{DH})}{[1-n(\mathrm{DH})]\times0.64^{0.5}}=1.0$，即 $n(\mathrm{DH})=0.8[1-n(\mathrm{DH})]$，$n(\mathrm{DH})=\dfrac49$ mol；吸收 HCl 的量 $=\beta=x\cdot n(\mathrm{DH})=0.5\times\dfrac49\approx0.22$ mol。
\item (1) 设生成 \ce{N2O4} 的分压为 $x$：$p(\ce{NO2})=100-2x$，总压 $100-x=70$，$x=30$ kPa；$p(\ce{NO2})=40$ kPa，$p(\ce{N2O4})=30$ kPa，$K_p=\dfrac{30}{40^2}=0.01875$ kPa$^{-1}$。(2) 平衡时 $\vz=\vn$，$K_p=\dfrac{\kz}{\kn}$。(3) 此时 $p(\ce{N2O4})=\dfrac{100-60}{2}=20$ kPa，$\dfrac{\vz}{\vn}=\dfrac{\kz\times60^2}{\kn\times20}=0.01875\times180=3.375$（$>1$，反应正向进行）。(4) 升温后 \ce{NO2} 的体积分数增大，说明平衡逆向移动，逆反应吸热，正反应放热，$\Delta H<0$。
\item (1) $<$；低温、高压。升温 $x$ 减小，平衡逆移，正反应放热；加压 $x$ 增大，有利于气体分子数减小的正反应。(2) 设起始 CO 为 1 mol，生成 \ce{Ni(CO)4} $y$ mol：剩余 CO $(1-4y)$ mol，气体总量 $(1-3y)$ mol。$\dfrac{y}{1-3y}=0.90$，$y=\dfrac{0.90}{3.7}\approx0.243$；$\alpha(\ce{CO})=4y\approx97.3\%$。$p(\ce{Ni(CO)4})=0.90$ MPa，$p(\ce{CO})=0.10$ MPa，$K_p=\dfrac{0.90}{0.10^4}=9.0\times10^3$ MPa$^{-3}$。(3) 低温时平衡正移，粗镍中的 Ni 转化为气态 \ce{Ni(CO)4} 与杂质分离；高温时平衡逆移，\ce{Ni(CO)4} 分解析出纯镍，CO 可循环使用——利用温度对平衡的影响实现“搬运”提纯。
\end{jiexilist}

